% ===================================================================
%  TÁBLA Uimh. 1 — An Scoil. — An Lycée. — An Rang.
%
%  Ceci n'est PAS une transcription : c'est une traduction, faite en
%  2026, en irlandais standard (An Caighdeán Oifigiúil). D'ou trois
%  differences avec texto/io et texto/fr, et trois seulement :
%
%   * pas de \nl ni de \cc ;
%   * pas de \begin{VUpage} ;
%   * le gras se pose sur le TERME seul, ponctuation dehors.
%
%  Les cles « %%K » sont celles de texto/io, macro pour macro pour
%  l'apparat, et les renvois \textsuperscript{(n)} visent les MEMES
%  objets de la planche.
%
%  ONZIEME COLONNE DES DIX-SEPT, et la premiere langue celtique de la
%  serie. Elle apporte deux traits que le volume n'avait pas encore
%  rencontres.
%
%  LE PREMIER EST L'ORDRE VSO : le verbe passe devant le sujet.
%  « Suíonn an múinteoir » — « s'assied le maitre ». Aucune autre
%  colonne des seize precedentes ne fait cela.
%
%  MAIS CELA NE COUTE AUCUNE INVERSION, et il faut le dire tout de
%  suite parce que c'est contre-intuitif : ce qui deplace un renvoi
%  n'est pas la place du verbe, c'est la place du complement du nom.
%  L'irlandais postpose le sien comme l'ido — « doras an chaisleáin »,
%  la porte du chateau — et l'ordre des SUBSTANTIFS entre eux ne bouge
%  donc pas. On attend zero inversion.
%
%  LE SECOND EST LA MUTATION INITIALE. Le mot change de commencement
%  selon ce qui le precede : « an chéad shraith », ou « céad » devient
%  « chéad » et « sraith » devient « shraith ». Le gras porte donc la
%  forme MUTEE, celle que la phrase demande, et le releve des objets
%  s'en trouve intitule comme la phrase le dit — meme parti qu'en
%  russe et en basque pour la declinaison.
%
%  CELA OBLIGE LE CABLAGE A ECRIRE LES DEUX ETATS : ORDINALO recoit
%  « ch?[eé]ad », SERIO recoit « sh?raith ». C'est la premiere fois de
%  la serie qu'un membre de ces groupes doit prevoir une alteration a
%  l'INTERIEUR du mot, et non seulement une desinence a la fin.
% ===================================================================

%%K t01-apar-0 apar
\VUsaut{2.0mm}
\VUcentre{12.6pt}{120}{LEABHRÁN MÍNIÚCHÁIN}
\VUsaut{3.2mm}
\VUcentre{8.4pt}{160}{DO}
\VUsaut{2.6mm}
\VUtitre{15.6pt}{0}{tháblaí cabhracha Delmas}
\VUsaut{3.4mm}
\VUornamento{9pt}
\VUsaut{5.0mm}
\VUcentre{13.2pt}{90}{AN CHÉAD SHRAITH}
\VUsaut{3.0mm}
\VUfilet{16mm}
\VUsaut{5.6mm}

%%K t01-apar-1 apar
\VUtitre{15.0pt}{60}{TÁBLA Uimh. 1}
\VUsaut{1.4mm}
\VUtitre{11.4pt}{0}{An Scoil. --- An Ardscoil. --- An Rang.}
\VUsaut{2.2mm}
\VUfilet{20mm}
\VUsaut{6.4mm}

%%K t01-tit-1 sub
\VUpk{10.2pt}{40}{Cur síos ginearálta.}
\VUsaut{3.0mm}

%%K t01-01-1 p
1. --- Sa tábla seo léirítear \VUgras{rang} in \VUgras{ardscoil}. Sa rang seo tá ocht
\VUgras{mbord} \textsuperscript{(18)} agus ocht \VUgras{mbinse}
\textsuperscript{(32)}. Ar gach binse suíonn triúr dalta. Suíonn an
\VUgras{múinteoir} \textsuperscript{(7)} ar \VUgras{chathaoir}
\textsuperscript{(8)} taobh thiar dá \VUgras{dheasc} \textsuperscript{(6)}.

%%K t01-02-1 p
2. --- Ar chlé na deisce tá \VUgras{cófra} gloine \textsuperscript{(15)}. Ar dheis
tá an \VUgras{clár dubh} \textsuperscript{(1)} lena \VUgras{cheirt}
\textsuperscript{(2)} agus lena \VUgras{chailc} \textsuperscript{(3)}.

%%K t01-02-2 p
Os cionn na deisce tá \VUgras{táblaí balla} \textsuperscript{(9, 11, 12)} agus
\VUgras{léarscáil} \textsuperscript{(10)} ar crochadh.

%%K t01-03-1 p
3. --- Níl na daltaí uile ina \VUgras{n-áiteanna} \textsuperscript{(17)}. Tá Eoin
\textsuperscript{(4)} ina sheasamh ag an gclár; tá an cheirt ina láimh aige agus tá
sé ag glanadh na \VUgras{bhfíoracha geoiméadracha}.

%%K t01-03-2 p
Tá Peadar ina sheasamh ag an deasc; tá sé ag bailiú na \VUgras{gcleachtaí
scríofa} \textsuperscript{(5)} chun iad a thabhairt don mhúinteoir.

%%K t01-04-1 p
4. --- Is múinteoir \VUgras{teangacha nua-aimseartha} é \VUgras{an} múinteoir seo.
Múineann sé do na daltaí conas labhairt, léamh agus scríobh i dteangacha
iasachta \VUgras{(Gearmáinis, Béarla, Spáinnis, Iodáilis, Rúisis} nó
\VUgras{Fraincis)}.

%%K t01-05-1 p
5. --- Tá an rang an-mhór agus an-solasmhar. Ag an gcúl tá \VUgras{fuinneog}
leathan le dhá \VUgras{scáthlán} \textsuperscript{(78)} agus \VUgras{doras
gloine}, chun é a aerú agus a shoilsiú.

%%K t01-05-2 p
Ní bhíonn fuacht sa rang seo. Sa lár, chun é a théamh, tá \VUgras{sorn}
an-mhaith \textsuperscript{(46)} lena \VUgras{phíopa} \textsuperscript{(45)}.

%%K t01-05-3 p
Sa landair tá \VUgras{crúcaí} \textsuperscript{(58)} buailte; orthu crochtar
hataí agus cótaí na ndaltaí.

%%K t01-tit-2 sub
\VUsaut{5.2mm}
\VUpk{10.2pt}{40}{An Chlós.}
\VUsaut{3.6mm}

%%K t01-06-1 p
6. --- Léiríonn an \VUgras{clog} \textsuperscript{(63)} cúig nóiméad tar éis a naoi.

%%K t01-06-2 p
Tá an \VUgras{sos} \textsuperscript{(76)} díreach thart, agus tá an ceacht ag tosú
arís. Bíonn an sos sa \VUgras{chlós} \textsuperscript{(75)}. Feicimid roinnt
daltaí ag imirt. Tá an \VUgras{maor} \textsuperscript{(74)} ag faire orthu.

%%K t01-07-1 p
7. --- Sa chlós feicimid daoine eile freisin: an \VUgras{cigire}
\textsuperscript{(73)}, an \VUgras{stiúrthóir} \textsuperscript{(71)} agus an
\VUgras{ceann staidéir} \textsuperscript{(72)}.

%%K t01-07-2 p
Scrúdaíonn an cigire na scoileanna, riarann an stiúrthóir an foras, agus faireann
an ceann staidéir ar obair na ndaltaí.

%%K t01-07-3 p
Is é an ceathrú duine an \VUgras{doirseoir} \textsuperscript{(77)}. Iompraíonn sé
an clár faoina ascaill, chun na daltaí atá as láthair a mharcáil.

%%K t01-08-1 p
8. --- Ag cúl an chlóis tá na \VUgras{fearais ghleacaíochta}
\textsuperscript{(70)} agus ceardlann na \VUgras{hoibre láimhe}
\textsuperscript{(69)}.

%%K t01-08-2 p
Ar dheis an tábla is ar éigean a fheictear \VUgras{ranganna} an \VUgras{cheoil}
agus na \VUgras{líníochta}.

%%K t01-noto-1 noto
\VUnotes{143.2mm}{%
(*) Do \textit{na hainmneacha} freisin molaimid iad a thabhairt sa fhoirm
Laidine. Sin an t-aon slí chun na hathruithe gan áireamh a sheachaint. Agus conas
a roghnófaí idir \textit{Giovanni, Jean, Johann, Jan, Hans, John, Eoin} agus a
thuilleadh? Ní dhéanfaimid foclóir ar leith d'ainmneacha. Glacaimis leis an
ainmneach Laidine agus abraimis: \VUgras{Ioannes, Ioanna; Iulius, Iulia;
Iakobus; Andreas; Lukas.} (Gramadach iomlán).%
}

%%K t01-tit-3 sub
\VUpk{10.2pt}{40}{Cur síos mion.}
\VUsaut{2.4mm}

%%K t01-tit-4 sub
\VUpk{10.2pt}{40}{Na daltaí maithe.}
\VUsaut{3.0mm}

%%K t01-09-1 p
9. --- Is iad na daltaí ar an gcéad bhinse ar dheis na deisce ná \VUgras{Anraí}
(*), \VUgras{Stiofán} agus \VUgras{Cathal}.

%%K t01-09-2 p
Tá Anraí an-aireach agus an-dícheallach; éisteann sé leis an múinteoir. Ar an
mbord aige tá \VUgras{cóipleabhar} \textsuperscript{(28)}, \VUgras{leabhar}
\textsuperscript{(29)}, \VUgras{foclóir} \textsuperscript{(30)} agus \VUgras{cás
peann luaidhe} \textsuperscript{(31)}. Ina leabhar tá go leor \VUgras{leathanach};
i ngach caibidil tá roinnt \VUgras{alt}, agus i ngach \VUgras{líne} tá go leor
\VUgras{focal}.

%%K t01-09-4 p
Tá a chomharsa Stiofán \textsuperscript{(24)} saothrach. Tá sé ag scríobh sa
chóipleabhar an cheachta atá tugtha don chéad uair eile. Tá \VUgras{lámh
scríbhneoireachta} den scoth aige. Tá a leabhar dúnta. Ní fheictear ach an
\VUgras{clúdach} \textsuperscript{(27)}. Lena thaobh tá a \VUgras{chompás}
\textsuperscript{(26)}.

%%K t01-09-5 p
Coinníonn Cathal \textsuperscript{(23)} an leabhar ina láimh; osclaíonn sé é chun
an ceacht a athrá uair amháin eile. Mar a bhíonn ag gach dalta, tá
\VUgras{dúchán} \textsuperscript{(25)} os a chomhair. Tá sé umhal.

%%K t01-10-1 p
10. --- Ag an gcéad bhord eile suíonn \VUgras{Labhrás, Antoine} agus
\VUgras{Pól}. Aithrisíonn Labhrás a \VUgras{cheacht} \textsuperscript{(22)} agus
freagraíonn sé \VUgras{ceisteanna} an mhúinteora. Tá Antoine
\textsuperscript{(21)} an-neamhaireach; uaireanta pionósaíonn an múinteoir é fiú. Is
\VUgras{comrádaí} maith é Pól \textsuperscript{(20)}, cairdiúil agus cabhrach. Tá
sé an-aireach.

%%K t01-tit-5 sub
\VUsaut{4.2mm}
\VUpk{10.2pt}{40}{Na daltaí dána.}
\VUsaut{3.2mm}

%%K t01-11-1 p
11. --- Taobh thiar d'Anraí suíonn \VUgras{Seoirse} \textsuperscript{(38)}. Ní
buachaill olc é, ach tá sé \VUgras{cainteach}, neamhaireach agus corrthónach.
Cuireann a \VUgras{éadromacht} agus a \VUgras{easumhlaíocht}
\VUgras{mí-ord} sa cheacht, agus is minic a thuilleann sé
\VUgras{rabhadh} géar. Tá a \VUgras{dhréacht} \textsuperscript{(35)} salach;
déanann sé \VUgras{smeadar dúigh} air lena \VUgras{pheann}
\textsuperscript{(34)}, agus dá bhrí sin bíonn a \VUgras{scriosán}
\textsuperscript{(36)} de shíor á úsáid aige; tá a \VUgras{mhála scoile}
\textsuperscript{(33)} stróicthe, chomh míshlachtmhar sin atá sé. Cad chuige a
dtugann sé a \VUgras{pheann tobair} \textsuperscript{(37)} go rang na dteangacha
nua-aimseartha?

%%K t01-11-2 p
Ní maith le Éamann \textsuperscript{(39)}, a chomharsa, é. Is fíor go bhfuil sé
beagán leisciúil, ach tá sé ciúin agus socair. Ní bhíonn sé ag caint; ach in ionad
éisteacht, is fearr leis \VUgras{scéalta} a léamh ina \VUgras{leabhar léitheoireachta}.

%%K t01-11-3 p
Is é an tríú dalta ar an mbinse seo ná \VUgras{Ludwig} \textsuperscript{(40)}. Tá
sé dícheallach. Scríobhann sé ar \VUgras{bhileog pháipéir} bhán. Os a chomhair chuir
sé an \VUgras{páipéar súite} \textsuperscript{(41)}.

%%K t01-12-1 p
12. --- Tá na daltaí ar an dara binse, ar chlé an mhúinteora, an-óg. An chéad
duine, \VUgras{Eimile} beag, níl sé in ann scríobh go maith fós; tugann sé leis a
\VUgras{shlinn} \textsuperscript{(42)}, a \VUgras{pheann slinne} agus a
\VUgras{spúinse} \textsuperscript{(43)}.

%%K t01-12-2 p
Lena thaobh tá \VUgras{Ágastas} agus \VUgras{Bearnard}
\textsuperscript{(44)}. Oibríonn an duine deireanach seo go maith, ach tá
\VUgras{cuma} an-mhíshlachtmhar air.

%%K t01-13-1 p
13. --- Taobh thiar díobh suíonn triúr \VUgras{leiscéirí}. Ní fhoghlaimíonn
\VUgras{Ailbe} a chuid ceachtanna. Bíonn \VUgras{Séamas} \textsuperscript{(47)} ina
chodladh in ionad éisteacht. Tarraingíonn a \VUgras{leisce}
\VUgras{iomardaithe} agus \VUgras{pionóis} fiú air. Faoi dheireadh, tá
\VUgras{Críostóir} \textsuperscript{(48)} róéadrom. \VUgras{Iompraíonn} sé é féin go
holc agus ní dhéanann sé iarracht feabhas a chur air féin.

%%K t01-14-1 p
14. --- A gcomharsana, os a choinne sin, iompraíonn siad iad féin go maith. Is
maith le \VUgras{Earnán} \textsuperscript{(51)} ord agus cruinneas; úsáideann sé
i gcónaí a \VUgras{riail} \textsuperscript{(50)} agus a \VUgras{pheann luaidhe}
\textsuperscript{(49)}. Is \VUgras{dalta lae} é.

%%K t01-14-2 p
Is \VUgras{cónaitheoir} é \VUgras{Proinsias} \textsuperscript{(52)}; caitheann sé
éide, itheann sé agus codlaíonn sé san fhoras. Is \VUgras{comrádaí} maith é.

%%K t01-14-3 p
Bíonn Muiris ag cur faobhair ar a \VUgras{pheann luaidhe} lena \VUgras{scian
phóca} \textsuperscript{(56)}; tá sé an-tíosach.

%%K t01-14-4 p
Tugann sé leis a chuid leabhar go léir, a \VUgras{ghramadach}
\textsuperscript{(55)}, a \VUgras{leabhrán nótaí} \textsuperscript{(54)} agus fiú a
\VUgras{cheap nótaí} \textsuperscript{(53)}. Tá a \VUgras{mhála cáipéisí}
\textsuperscript{(57)} ar an urlár taobh thiar de.

%%K t01-14-5 p
Tá an dá bhord ag cúl an ranga ag na \VUgras{daltaí maithe}
\textsuperscript{(19)}.

%%K t01-tit-6 sub
\VUsaut{4.6mm}
\VUpk{10.2pt}{40}{An líníocht agus an ceol.}
\VUsaut{3.4mm}

%%K t01-15-1 p
15. --- I \VUgras{rang na líníochta} tá \VUgras{samplaí} áille ar crochadh ar an
mballa, agus ón tsíleáil crochtar \VUgras{lampa} \textsuperscript{(62)} lena
\VUgras{scáthlán}. Tá an \VUgras{múinteoir} \textsuperscript{(60)} an-fhoighneach;
ceartaíonn sé \VUgras{líníochtaí} \textsuperscript{(61)} na ndaltaí agus múineann
sé dóibh conas an \VUgras{busta} \textsuperscript{(59)} a tharraingt ón nádúr.

%%K t01-16-1 p
16. --- Dealraíonn \VUgras{rang an cheoil} an-taitneamhach. Ann foghlaimítear conas
na \VUgras{nótaí} a léamh, conas \VUgras{céimeanna an scála}
\textsuperscript{(67)} a chanadh --- scríofa ar an \VUgras{chláirseach cheoil} (do,
re, mi, fa, sol, la, si\,=\,C. D. E. F. G. A. B.) ---, conas na \VUgras{eochracha}
a aithint agus \VUgras{foinn} áille a chanadh. Cuirtear na nótaí ar an
\VUgras{seastán} \textsuperscript{(65)}. Bíonn duine de na daltaí \VUgras{ag seinm
an phianó} \textsuperscript{(64)}, agus bíonn \VUgras{múinteoir an cheoil}
\textsuperscript{(66)} \VUgras{ag bualadh an ama} \textsuperscript{(68)}.

%%K t01-17-1 p
17. --- Is ar éigean a fheictear cúinne de cheardlann na \VUgras{hoibre láimhe}
\textsuperscript{(69)}, mar a bhféadann na páistí foghlaim conas adhmad a
locairt agus iarann a líomhadh. Ní bhíonn a leithéid i ngach scoil.

%%K t01-tit-7 sub
\VUsaut{5.0mm}
\VUpk{10.2pt}{40}{Rang na n-eolaíochtaí nádúrtha.}
\VUsaut{3.6mm}

%%K t01-18-1 p
18. --- Múineann múinteoir na matamaitice do na daltaí \VUgras{céimseata,}
\VUgras{uimhríocht, áireamh} agus an \VUgras{córas méadrach}. Sa tábla feicimid na
príomhfhíoracha céimseata: an \VUgras{ciorcal} \textsuperscript{(a)}, an
\VUgras{cearnóg} \textsuperscript{(b)}, an \VUgras{triantán}
\textsuperscript{(c)}, an \VUgras{líne} \textsuperscript{(d)}.

%%K t01-18-2 p
Feicimid \VUgras{oibríocht} freisin; ní \VUgras{suimiú} é, ní \VUgras{dealú} é
ná \VUgras{iolrú}: is \VUgras{roinnt} \textsuperscript{(e)} é.

%%K t01-19-1 p
19. --- I dtábla an chórais mhéadraigh feicimid: an \VUgras{méadar}
\textsuperscript{(a)}, an \VUgras{scála} \textsuperscript{(b)} agus a chuid
\VUgras{meáchan}, an \VUgras{lítear} \textsuperscript{(e)}, an
\VUgras{deacailítear} \textsuperscript{(f)}, an \VUgras{deicilítear} agus an
\VUgras{méadar ciúbach} \textsuperscript{(d)}.

%%K t01-19-2 p
Foghlaimíonn na daltaí na \VUgras{heolaíochtaí nádúrtha} freisin
\textsuperscript{(11)} --- zó-eolaíocht \textsuperscript{(a)}, luibheolaíocht
\textsuperscript{(b)}, geolaíocht \textsuperscript{(c)} --- agus an
\VUgras{stair} \textsuperscript{(12)}.

%%K t01-19-3 p
Sa chófra tá gléasanna don \VUgras{fhisic} \textsuperscript{(15)} agus don
\VUgras{cheimic} \textsuperscript{(16)}.

%%K t01-19-4 p
Ar an gcófra tá: an \VUgras{sféar} \textsuperscript{(14)}, an \VUgras{cón} agus an
\VUgras{cruinneog} \textsuperscript{(13)}.

%%K t01-19-5 p
Os cionn na dtáblaí tá \VUgras{léarscáil} ar crochadh a léiríonn cúig chuid an
domhain: an \VUgras{Eoraip} \textsuperscript{(g)}, an \VUgras{Áise}
\textsuperscript{(e)}, an \VUgras{Afraic} \textsuperscript{(f)}, \VUgras{Meiriceá}
\textsuperscript{(ab)} agus an \VUgras{Aigéine} \textsuperscript{(h)}, agus an dá
aigéan mhóra --- an \VUgras{tAigéan Ciúin} \textsuperscript{(c)} agus an
\VUgras{tAigéan Atlantach} \textsuperscript{(d)}.
\VUsaut{6.0mm}
\VUfilet{22mm}
